%%%% CAPÍTULO 1 - INTRODUÇÃO
%%
%% Deve apresentar uma visão global da pesquisa, incluindo: breve histórico, importância e justificativa da escolha do tema,
%% delimitações do assunto, formulação de hipóteses e objetivos da pesquisa e estrutura do trabalho.

%% Título e rótulo de capítulo (rótulos não devem conter caracteres especiais, acentuados ou cedilha)
\chapter{Introdução}\label{cap:introducao}

%Neste capítulo, será apresentada uma perspectiva geral do trabalho a ser desenvolvido, demonstrando o contexto de aplicação e a relevância do desenvolvimento.

%\section{Considerações iniciais}\label{sec:consideracoesIniciais}

%Desde os primórdios, quando a humanidade iniciou o desenvolvimento de técnicas e o aprimoramento de suas criações, o impulso de superar limitações e aflições da realidade tem se configurado como a principal motivação para a evolução tecnológica. Observa-se que, desde a descoberta do fogo e o domínio da agricultura até as transformações proporcionadas pela Revolução Industrial e, mais recentemente, pela era digital, cada avanço histórico esteve intimamente relacionado ao desejo humano de transcender as restrições impostas pela natureza. Tal movimento evidencia que o desenvolvimento tecnológico não se origina exclusivamente da necessidade prática, mas também de uma inquietação existencial, caracterizada pela busca incessante por conforto, segurança e eficiência. O aprimoramento contínuo do que já foi concebido revela-se, portanto, como um fator essencial para o progresso constante da humanidade e para a construção de novos horizontes sociais, culturais e científicos \cite{tec:civ}.

%A necessidade intrínseca do ser humano de criar e aperfeiçoar aquilo que já foi desenvolvido propiciou, ao longo da história, mudanças radicais em diferentes esferas da realidade. Um exemplo paradigmático é o conhecimento matemático: inicialmente empregado para contagem e controle de recursos, tornou-se gradualmente fundamental para o planejamento e a construção de edificações, para a navegação e o mapeamento de territórios e, mais recentemente, para o desenvolvimento de computadores e sistemas digitais. O constante aprimoramento desse conhecimento evidencia como a evolução das técnicas e das ciências expande progressivamente os limites do que a humanidade é capaz de realizar.

No campo da computação, a evolução do hardware acompanha décadas de avanço: de válvulas a vácuo a circuitos de escala muito reduzida, até os computadores atuais \cite{eniac}. Essa trajetória levou à Internet e a novas áreas de estudo, com busca recorrente por mais eficiência, desempenho e confiabilidade nos sistemas.

No desenvolvimento de software, vale o mesmo princípio: muitos sistemas precisam de atualização contínua, refatoração do código e novas funcionalidades para não ficarem obsoletos. O desenvolvimento de software ocupa lugar central nesse cenário, por sustentar boa parte das soluções digitais usadas na educação, na cultura e na economia.

Nesse contexto de constante evolução tecnológica e de aprimoramento de sistemas, iniciativas educacionais e culturais que envolvem tecnologia digital demandam soluções adaptáveis e capazes de integrar novas funcionalidades de forma contínua. Reconhecendo essa necessidade de adaptação e aprimoramento, as universidades passaram a buscar maneiras de organizar e preservar informações de forma mais estruturada. A \acrfull{utfpr} e a \acrfull{unicentro}, a partir dos projetos de extensão "Tecno-Lixo: Oficina do Aprender" \cite{ribeiro2024tecno} e "E-Lixo" \cite{re2021cartilha}, respectivamente, que trabalham com o reaproveitamento de peças de computadores descartadas, acabaram armazenando diversos itens de informática em suas dependências para criação de um museu, preservando assim um pouco da história da computação. \citeonline{gonzaga2024} implementou um museu virtual de informática para cadastro e gerência desses itens, chamado E-Museu.

O E-Museu trata-se de um sistema web desenvolvido como espaço de catalogação, gerenciamento e exposição de itens tecnológicos no formato de museu digital. Esse sistema está em uso e encontra-se disponível na Web em \url{https://e-museu.gp.utfpr.edu.br}. Contudo, precisava de refatorações e de aprimoramentos além da adição de novas funcionalidades, motivando o presente trabalho. 
%Além da implementação de novas funcionalidades, tornou-se fundamental promover refatorações no código já implementado.% 
Dessa forma, o E-Museu ilustra essa lógica de evolução contínua e mostra como melhorias técnicas podem ampliar a experiência de quem consulta o acervo e de quem o administra.


Esse tipo de evolução enfrenta obstáculos: atualizar tecnologias sem quebrar o que já funciona, refatorar sem mudar o comportamento visível ao usuário e publicar o sistema em produção com confiabilidade. Este trabalho refatorou e ampliou o E-Museu com padrões de desenvolvimento de software, testes automatizados e ambientes padronizados, buscando tornar o sistema mais robusto, seguro e de fácil manutenção.



\section{Objetivos}\label{sec:objetivos}

\subsection{Objetivo geral}\label{subsec:objetivoGeral}

% Não coloque o padrão de código que vai ser usado por não ter essa noção ainda, acredito que da pra ver isso no futuro
Refatorar e implementar novas funcionalidades no E-Museu, aplicando boas práticas de programação, padrões de projeto de software e testes automatizados, com o intuito de melhorar a manutenibilidade e a qualidade do código.
%Implementar novas funcionalidades e refatorar o sistema Web do E-museu aplicando boas práticas de programação, padrões de projeto de software e testes.

%Refatorar e implementar novas funcionalidades no sistema web E-Museu pela UTFPR e UNICENTRO, visando aprimorar sua estrutura de código, aumentar a eficiência e a manutenibilidade, além de ampliar os recursos disponíveis para proporcionar uma experiência mais intuitiva, robusta e alinhada às demandas tecnológicas contemporâneas. <-- aqui já está quase justificando

\subsection{Objetivos Específicos}\label{subsec:objetivoEspecificos}

\begin{itemize}
    \item Estabelecer uma estrutura de versionamento e desenvolvimento de código com padrões bem definidos, utilizando GitHub, de modo a garantir organização, rastreabilidade e colaboração eficiente.
    \item Configurar ambientes estáveis de desenvolvimento, teste e produção, assegurando consistência, reprodutibilidade e confiabilidade do sistema.
    \item Refatorar o código existente, aprimorando legibilidade, modularidade, manutenibilidade e desempenho. %, além de reestilizar a interface para proporcionar uma experiência mais clara, intuitiva e acessível aos usuários.
    \item Implementar testes automatizados a fim de assegurar a qualidade, confiabilidade e robustez do código, contemplando testes de fluxo e de integração.
    %\item Desenvolver e aprimorar funcionalidades do E-Museu, como cadastro e gerenciamento de itens, upload de fotos e vídeos, geração de QR Code, relatórios de acesso e itens cadastrados, internacionalização do sistema e criação de newsletter, incorporando ainda mecanismos de segurança como reCAPTCHA e proteção de dados.
    \item Desenvolver e aprimorar funcionalidades do E-Museu, incluindo múltiplas imagens por item, internacionalização do sistema, proteção de formulários (por exemplo Turnstile) e geração de etiquetas com código e nome para as peças físicas.
    \item Aplicar padrões de projeto e boas práticas de desenvolvimento, orientando tanto a refatoração quanto a implementação de novas funcionalidades de forma estruturada e eficiente.
\end{itemize}

\section{Justificativa}\label{sec:justificativa}

%O E-Museu, em seu estado inicial, apresentava uma infraestrutura bastante simplificada, sem a separação entre os ambientes de desenvolvimento e produção — etapa fundamental no ciclo de vida de um software, quando o sistema atinge sua versão final e passa a operar no mundo real. Cada execução da aplicação demandava configurações manuais complexas, como a instalação de tecnologias necessárias ou a preparação do banco de dados. Além disso, o código não seguia práticas consolidadas de desenvolvimento, como a adequada divisão de responsabilidades entre \textit{services}, \textit{actions} e \textit{controllers} — elementos essenciais para a organização e manutenção de sistemas robustos. \cite{clean:arch}

Na versão analisada no início deste trabalho, o E-Museu não separava de forma clara os ambientes de desenvolvimento e produção na operação diária. Cada execução da aplicação exigia configurações manuais (instalação de dependências \acrshort{php}, preparação do banco de dados e do servidor web). O código seguia uma versão simplificada do padrão \acrfull{mvc}, com lógica concentrada em dezoito controladores, sem divisão consistente entre \textit{services}, \textit{actions} e controladores, o que dificultava a manutenção.

% Além disso, o código não utiliza padrões de projeto e segue práticas consolidadas de desenvolvimento, como a adequada divisão de responsabilidades entre \textit{services}, \textit{actions} e \textit{controllers} — elementos essenciais para a organização e manutenção de sistemas robustos. 
Faltavam também mecanismos automatizados de qualidade e de validação do comportamento da aplicação, como \textit{linters} e testes automatizados. As funcionalidades existentes precisavam de evolução, o que reforçou a necessidade de refatoração e de novos recursos \cite{importance:testing}. %, escalável e alinhada às demandas atuais 

%Outro ponto crítico é a ausência de testes automatizados, tanto de qualidade quanto de funcionamento, os quais atualmente são indispensáveis para garantir a confiabilidade e a evolução de qualquer aplicação. Somado a isso, as funcionalidades existentes podem ser melhoradas, as quais reforçam a necessidade de refatorações e da elaboração de novos recursos, de modo a tornar a plataforma mais consistente, escalável e alinhada às demandas atuais. \cite{importance:testing}

Nesse contexto, a refatoração assume papel central, pois melhora a estrutura interna do código sem alterar o comportamento visível ao usuário, trazendo mais clareza, modularidade e facilidade de manutenção. Ao reduzir complexidade e duplicações, a refatoração facilita correções e novas funcionalidades com menor risco de regressão \cite{refatoracao:2004}. %Paralelamente, a restilização da interface, ajustando cores, tipografia e detalhes visuais, assume relevância estratégica, uma vez que influência diretamente a usabilidade e a percepção dos usuários, favorecendo uma experiência mais intuitiva, eficiente e satisfatória no uso do sistema.

\section{Estrutura do trabalho}\label{sec:estruturaTrabalho}

O presente documento está organizado em sete capítulos, além dos elementos pré-textuais e pós-textuais exigidos pelo template da UTFPR.

O Capítulo~\ref{cap:introducao} apresenta o contexto, a justificativa, os objetivos e a estrutura desta monografia.

O Capítulo~\ref{cap:trabalhos-relacionados} relaciona o trabalho a versões anteriores do E-Museu e a produções acadêmicas ou sistemas comparáveis em museologia digital e engenharia de software.

O Capítulo~\ref{cap-referencial-teorico} reúne os fundamentos teóricos sobre o E-Museu, refatoração de software e boas práticas de desenvolvimento adotadas no projeto.

O Capítulo~\ref{cap:materiais-e-metodos} descreve os materiais (tecnologias, ferramentas e infraestrutura) e os métodos empregados no planejamento e na execução do trabalho.

O Capítulo~\ref{cap:desenvolvimento}, núcleo da monografia, documenta o desenvolvimento realizado: requisitos, versionamento, ambientes, refatoração, testes, funcionalidades e implantação do sistema.

O Capítulo~\ref{cap:resultados} consolida os resultados alcançados: síntese frente aos objetivos, evolução do modelo de dados e do software em relação ao legado, indicadores do repositório e entrega operacional.

Por fim, o Capítulo~\ref{cap:consideracoesfinais} retoma os objetivos à luz do que foi realizado, explicita limitações e indica trabalhos futuros para a continuidade do E-Museu, inclusive como base para futuros Trabalhos de Conclusão de Curso sobre o mesmo sistema.
